The transition towards converter-dominated power systems, driven by the massive integration of renewable energy sources, is fundamentally altering the dynamic behaviour of electrical grids. As synchronous generation is progressively displaced by inverter-based resources, new stability challenges emerge: reduced system inertia, faster electromagnetic transients, and complex interactions between converter controls and the network. In this context, transmission system operators require tools that can assess stability rapidly, reliably, and under a wide range of operating scenarios.

The European project \textit{TenSyGrid} addresses this need by developing a toolbox for direct stability assessment based on a tensor (multilinear) reformulation of power system component models. By expressing the nonlinear dynamics of generators, converters, lines, loads, and transformers in a multilinear form, the toolbox enables efficient computation of the system's state matrix and its eigenvalues, which are the fundamental indicators of small-signal stability. The toolbox is designed to be compatible with existing commercial software packages, facilitating its integration into current operator workflows.

This deliverable (D4.1) constitutes the first step in Work Package~4 (\textit{TenSyGrid Toolbox Development}). Its objective is to identify, describe, and catalogue the use cases that the toolbox must support, taking the network operator as the primary user. The use cases are derived from three complementary sources:

\begin{enumerate}
  \item The TenSyGrid project objectives, as defined in the CETPartnership proposal.
  \item An initial web-based \textit{Stability Monitor} prototype that implements core stability-analysis workflows.
  \item An anticipated operator-needs analysis, to be validated during a dedicated workshop with grid operators.
\end{enumerate}

The document is structured as follows. Chapter~2 describes the methodological approach, the evolution of the software architecture towards a native open-source Python framework, and the insights gathered from operator interviews. Chapter~3 presents the use case inventory, organised into implemented and proposed use cases, and outlines the organization of the toolbox interface. Each use case is described in terms of its actor, objective, theoretical basis, interaction flow, and expected outcome. The chapter also documents use cases that were considered but discarded or deferred, together with the rationale for their exclusion. Chapter~4 concludes the deliverable and outlines the next steps towards the operator workshop and the subsequent development cycle.

The intended audience for this deliverable includes the TenSyGrid consortium partners, the CETPartnership review bodies, and --- after the operator workshop --- the network operators who will ultimately use and benefit from the toolbox.
